\subsection{Summary and Outlook}

This study addressed the gap between detecting abnormal UAV telemetry and producing diagnostic evidence that can guide engineering inspection. HS-AeroTS organizes this process as a reproducible hierarchy: a first-stage detector identifies runtime anomalies, a second-stage classifier assigns anomalous windows to four fault domains, TreeSHAP contributions are aggregated into traceable channel, topic, and subsystem evidence, and a static PX4 uORB graph propagates matching-commit topic evidence into software-module suspect rankings. The framework therefore connects predictive decisions with architecture-aware diagnostic evidence while preserving the distinction between runtime annotations and source-code faults.

The experiments support the utility of the hierarchy within this defined scope. On the fixed chronological UAV-SEAD protocol, Stage 1 achieved an AUPRC of $0.7522 \pm 0.0039$, and the anomaly-only Stage 2 classifier achieved a Macro-F1 of $0.8980 \pm 0.0020$. SHAP aggregation conserved evidence mass and SHAP-ranked masking caused substantially greater performance degradation than random masking, but semantic-map agreement was sensitive to the mapping definition and weak for External and Global Position. The architecture graph covered all 18 evaluated topics and connected them to 54 PX4 modules. Its aggregate ranking was recomputed only from the 1,393 anomalous test windows that matched the selected commit; parser and degree sensitivities nevertheless bound the ranking to an inspection-support role.

The strict evaluations delimit these positive findings. Purged and leave-log-out isolation reduced complete-task performance relative to the fixed chronological protocol, and the leave-log-out comparison did not support a stable Macro-F1 difference between the Cascade and Direct Five-Class formulations. More importantly, the unchanged real-flight detector identified none of the 12 controlled mutation runs. Consequently, the nonzero onset-conditioned module ranks cannot be interpreted as detector-gated end-to-end localization performance. Static topic-to-module propagation provides version-dependent inspection evidence, but it neither proves fault causality nor validates an end-to-end module diagnosis capability.

Future work should prioritize a larger mutation benchmark spanning independent source logs, PX4 versions, software modules, and fault mechanisms, with development and test mutations kept separate. Such a benchmark would permit detector adaptation to be evaluated without contaminating real-flight or mutation test sets. Combining the current static uORB evidence with runtime topic activity, execution traces, and temporal propagation may further improve suspect ranking, but these extensions should continue to be assessed against known injected modules and under detector-gated protocols. These steps are necessary to determine whether the present architecture-aware evidence chain can advance from a traceable inspection aid toward dependable module-level diagnosis.
